\section{Conclusions}

Air--surface--underwater cooperation turns trust into an end-to-end mission problem. A defensible decision integrates a plausible observation, an intact and accountable provenance chain, a correctly governed shared state where one is needed, and evidence that remains suitable for the particular mission consequence. The authors' OAL framework connects these four assurance layers to entity, platform/device, link/path, and data-product objects over the complete lifecycle from admission to recovery.

The proposed object-dependency graph and CTG loop organize evidence flow, shared failure causes, uncertainty, and policy action across that structure. They distinguish competing explanations of system behavior from insufficient evidence and require numerical assessments to state their meaning and support. The evidence profile identifies complementary forms of analysis, implementation, environmental testing, and operational experience without treating them as a cumulative quality score.

Blockchain is a candidate coordination mechanism where independent organizations need shared membership and version state, attestation and provenance commitments, accountable audit, and post-partition reconciliation. Its value depends on governance, fault assumptions, and measured costs relative to signed logs and replicated databases. Physical corroboration supports observation validity; attestation appraises measured endpoint state; local safety mechanisms address communication outages and deadlines; and privacy and governance require additional controls. None of these component guarantees alone establishes mission suitability.

The proposed CDUTP and six research hypotheses define work toward interoperable, testable systems. The contribution of this review is the structure of the assurance questions and their connection to selected evidence; performance benefits remain to be established. The next stage combines formal analysis, cross-domain testbeds, representative field trials, and multi-party governance exercises, measuring both decision consequences and the quality of the supporting evidence. The practical objective is to make each consequential mission use traceable to current, appropriately scoped assurance and explicit residual uncertainty.
